% before: hf-frame
The host fields of this chapter are abelian, and their Galois groups all have the same shape: two
signs, one for each of the quadratic fields $\mathbf Q(\sqrt{-q})$ and $\mathbf Q(\sqrt{-d})$, and
a finite abelian group that maps onto the Galois group of a layer of the $2$-power cyclotomic tower.
Fixing this shape in advance lets the later chapters index places, characters and Frobenius classes
by elements of one explicit group.

% before: hf-host-field-of-character
The bound on Tate-Shafarevich groups in Chapter~\ref{chapter-tate-shafarevich} is proved for a
field described by Galois theory alone, as a subfield of a cyclotomic field. That description uses
only $\psi$, the exponent $e$ and the two primes, which is all that is fixed when the bound is
proved.

% before: hf-framed-host-field
The constructions on the definite side need more than the field: an identification of its Galois
group with the group of a Galois datum, and an Artin map from residues modulo the conductor that is
compatible with that identification. The conditions on the Artin map say which signs and which
part of the $2$-power tower each residue class controls.

% before: hf-host-field-properties
The data of a host field determine its arithmetic. The CM structure, the Frobenius elements and the
ramification at $d$ and in $E/F$ are read off from the cyclotomic field that contains $E$. The
ideal $\mathfrak b$, a square root of the relative different of $E/K$, makes $\mathfrak b^{-1}$
self-dual for the form $\operatorname{Tr}_{E/K}(y\bar y')$, which the Hermitian lattices of
Chapter~\ref{chapter-unitary-theta} use.

% before: hf-kernel-host-field
A host field built from $\psi$ is specified by the kernel of its Artin map. When that kernel is the
one cut out by $\psi$, the $2$-power tower and the two quadratic characters, the field is a host
field of $(\psi,e,q,d)$ in the sense used for the Tate-Shafarevich bound.

% before: hf-host-field-rows
A few further arithmetic conditions on a host field are needed for the definite side and for the
choice of auxiliary primes. They concern the quadratic subfield $\mathbf Q(\sqrt{-q})$, the
minimality of $Q_S$, and the splitting in $\mathbf Q(\sqrt{-d})$ of all small odd primes.

% before: hf-host-field-exists
For a fixed primitive character $\psi$ the host fields form an explicit family indexed by the
exponent $e$ and the pair $(q,d)$. The degree grows like $2^e$ with a constant factor that depends
only on $\psi$, which is what allows the degree to be matched to the scale $x$.

% before: hf-admissible-pair
The auxiliary primes have size $\exp(Y^{\varepsilon/5})$ and $\exp(Y^{\varepsilon/5}\log^2Y)$. The
congruence conditions make every small odd prime split in $\mathbf Q(\sqrt{-d})$ and fix the
behaviour of $2$, $q$ and $d$ in the quadratic fields involved.

% before: hf-good-pair
The characters of a host field built from $\psi$ are products of a power of $\psi$ with the two
quadratic characters attached to $q$ and $d$. Those that involve $q$ or $d$ must have square-root
cancellation at the scales of the window, and the real ones must have no zeros near $1$; this
definition states both requirements for every modulus $M$ at once.

% before: hf-zero-density-box
Most primes $q$ must give twisted characters without zeros near $s=1$. The zero-density estimate
bounds the number of exceptions, because a character with conductor $fq$ and $f$ bounded
determines $q$.

% before: hf-zero-density-regions
The definite side uses the classical zero-free regions together with a count of the zeros near
$s=1$ with multiplicity. They are stated here with the other analytic input on Dirichlet
$L$-functions and applied in Chapter~\ref{chapter-assembly}.

% before: hf-primes-progressions
The mean-square bound for twisted sums rests on the distribution of primes in progressions to
moduli that grow with $Y$. Only one real character can have a zero very close to $1$ among all the
moduli involved, and the count of primes in a progression depends on that zero through one explicit
term.

% before: hf-twisted-prime-sums
Expanding the square of a twisted sum leads to sums of Jacobi symbols $\left(\frac qs\right)$ over
the primes $q$ of a progression. These cancel for every modulus $s$ except possibly one, the one
attached to an exceptional zero, and the cancellation is by an arbitrary power of $Y$.

% before: hf-prime-power-sums
The diagonal terms of the expanded square are sums of the weights $\Lambda(n)\min(1,y/n)$ over
prime powers. Their size is that of a single sum of length about $y$, up to a logarithm, and this
gives the factor $y\log y$ in the mean square.

% before: hf-mean-square
Square-root cancellation for most twists comes from a mean value over the primes $q$ of a
progression. The twisted sums are truncated at a large power of $Y$, and the mean square is of the
size $y\log y$ that square-root cancellation predicts for each term.

% before: hf-log-derivative-box
A zero-free box near $s=1$ bounds the logarithmic derivative of $L(s,\chi)$ by $\log Q$ on a
slightly smaller box. This bound is used on the horizontal sides and on the left side of the contour
in the next lemma.

% before: hf-box-tail
The tail of the series $\sum_n\Lambda(n)\chi(n)/n$ beyond a large power of $Y$ is controlled by a
difference of two smoothed sums. Shifting the contour of Perron's formula into the zero-free box
bounds that difference, with a loss of only $z^{-1/400+\delta}$ relative to its trivial size.

% before: hf-chebyshev-from-twisted-sum
The mean-square bound controls truncated sums, while the window bound needs the corrected Chebyshev
sum $g_\chi$. A zero-free box near $s=1$ controls the tail of the series and the value
$\frac{L'}{L}(1,\chi)$, so a small truncated sum gives a small $g_\chi$.

% before: hf-most-primes-good
The two primes are chosen one after the other, $q$ first and then $d$ subject to congruences
modulo $q$ and the small odd primes. In both steps all but a small proportion of the primes in the
progression are good.

% before: hf-good-pairs-exist
Choosing $q$ and then $d$ among the good primes gives a pair that is good for every character
modulo $M$ at once. Here $M$ is the conductor of $\psi$ when the pair is used.

% before: hf-bounded-characters
In a host field built from $\psi$, a character whose primitive character is not a power of $\psi$
involves $q$ or $d$. A good pair therefore gives the bounds required for all these characters,
which are condition (e) of the host conditions.

% before: hf-windows
The window bound of Chapter~\ref{chapter-inert-prime-balance} is proved at scales $x$, and the
windowed form of Landau's theorem in Chapter~\ref{chapter-windowed-landau} needs it on a set of
scales that contains long intervals $[Y,Y^\lambda]$ arbitrarily far out. Each scale carries the
parameter $\varepsilon$ and an exponent $J$ that bounds how small $Y$ may be relative to $x$.

% before: hf-windows-exist
The bound on Tate-Shafarevich groups holds only on windows $[W,W^2]$ that depend on the pair
$(q,d)$, and the pair must be chosen before the window is known. Shrinking each such window to
exponent $3/2$ leaves room to rescale $x$ by the constant that relates $n$ to $2^e$.

% before: hf-no-short-relations
The thick places of a host are indexed by a set of integers placed along the cyclic subgroup $A_0$.
A Frobenius orbit of places moves along an arithmetic progression in $A$, and the condition below
prevents such a progression from passing through three thick places in a short stretch.

% before: hf-behrend-set
Behrend's construction gives sets without three-term progressions of density
$\exp(-C\sqrt{\log n})$. The same construction, read in a base large compared with $L$, excludes all
short linear relations at the cost of a factor depending on $L$.

% before: hf-behrend-set-at-scale
At a scale $x$ the set must be large compared with $x^{-\varepsilon}n$ while the length $L$ grows
like $\log x$. Behrend's bound loses only $\exp(O(\sqrt{\log x\log\log x}))$, which is smaller than any
power of $x$.

% before: hf-thick-orbits
The no-relation property is used through the following consequence. Along the orbit of a place
under the powers of a Frobenius element, the thick places occur at most twice in every stretch of length $L$ when the
orbit is long, and at most once when it is short.

% before: hf-adapted-ranks
The ranks of the unitary groups of Chapter~\ref{chapter-unitary-theta} are tied to the scale. The
order $n$ of the $2$-power part of the Galois group is about $x^{2+8\varepsilon}$, the number $m$
of theta lifts about $\sqrt x$, and the first rank $k$ about $x^{9/2+10\varepsilon}$; the twist
parameter $t$ is then determined by one equation.

% before: hf-adapted-ranks-exist
Ranks adapted to a scale exist for every Galois datum of the right order, and the twist parameter is
bounded by $x^{400}$. The second fact is why the Tate-Shafarevich bound is needed for twists up to a
fixed power of $x$.

% before: hf-host
A host collects the host field and the further choices used on the definite side: level primes at
which the theta kernels are twisted, a set of top classes and a Behrend set of thick places. The
conditions here are those that hold for every host; the estimates that depend on the scale come in
the definition of a good host.

% before: hf-host-extension
The data of a host other than the host field can be chosen freely subject to their stated
conditions. So a host is built by first constructing the host field and then choosing primes and
sets.

% before: hf-host-conditions
The comparison inequality of Chapter~\ref{chapter-inert-prime-balance} rewrites the weighted
Chebyshev sum of a level in terms of the characters of a host field. Its error terms are controlled
when the level's characters appear in $X^+$, the top classes are equidistributed over the cosets of
$A_0$ with the right weights, and all other characters have square-root cancellation at $x$.

% before: hf-good-orbit
The divisibility estimates of Chapter~\ref{chapter-period-divisibility} follow the top classes along
the orbits of Frobenius elements of inert primes. They need the positions of the top classes along
such an orbit to be sparse in every window, and every stretch of about $\log x$ consecutive tiles to
contain a tile in which single even positions alternate with empty cells.

% before: hf-good-host
The definite side uses many quantitative properties of a host at a given scale. They concern the
sizes of the auxiliary primes, analytic bounds for the quadratic characters and the real characters
of the host field, the norm primes $g_i$ and the other level primes, and the distribution of the top
classes and of the thick places.

% before: hf-norm-primes-short
The level primes $g_0,g_1$ must be norms of elements $\pi\equiv1$ of $\mathbf Q(\sqrt{-d})$ in a
dyadic interval. The number of such primes is bounded below in terms of the discriminant and the
modulus, with polynomial dependence on both.

% before: hf-exceptional-primes
Two norm primes are needed whose sets of exceptional primes are disjoint. A fixed prime is
exceptional for few integers of a dyadic interval, because it must divide $g^i-1$ to a high power, so
most pairs of norm primes have this property.

% before: hf-norm-primes
The level primes $g_0,g_1$ are norms of elements $\pi\equiv1$ from $\mathbf Q(\sqrt{-d})$, of size
about $(kb)^8$, with disjoint sets of exceptional primes. The other level primes are small primes
congruent to $1$ modulo the conductor of the host field.

% before: hf-pnt-character
At the scales $y\geq\exp(x^{\varepsilon/3})$ every character of a host field needs a Chebyshev
bound with a saving $\exp(-cx^{\varepsilon/8})$. The classical form of the prime number theorem for
a primitive character gives this saving once a real zero near $1$ is excluded.

% before: hf-chebyshev-integral
The corrected Chebyshev sum $g_\chi$ is an integral of the error term in the prime number theorem
for $\chi$. This identity transfers bounds for $\psi(u,\chi)-\delta_\chi u$ at all $u\geq y$ to a
bound for $g_\chi(y)$.

% before: hf-host-chebyshev
Every character of a host field, including the trivial one, must satisfy a Chebyshev bound at the
very large scales $y\geq\exp(x^{\varepsilon/3})$. At such scales the classical zero-free region
suffices once exceptional zeros are excluded.

% before: hf-orbit-failure
The top classes are chosen at random, and their positions along a long Frobenius orbit then form a
random set with independent entries. Each of the three conditions of a tiled set of positions fails
with probability at most a constant times $x^{-2}$.

% before: hf-random-selection
A random choice must satisfy two conditions at once: the distribution over the cosets of $A_0$,
controlled by second moments, and the weighted count of bad orbits, controlled by first moments. The
first fails with probability at most $1/2$ and the second with probability at most $1/4$, so some
choice satisfies both.

% before: hf-top-classes
The top classes must be equidistributed over the cosets of $A_0$ and must meet the Frobenius orbits
of the inert primes in a controlled way. A random choice with the right densities does both.

% before: hf-good-host-exists
Given a host field with a large splitting bound and a good pair at a scale, the remaining choices
of a good host can be made one at a time. The level primes, the Behrend set and the top classes are
chosen by the preceding lemmas, and the analytic conditions come from the good pair.

% before: hf-hosts-on-windows-def
Chapters~\ref{chapter-windowed-landau} and~\ref{chapter-assembly} need hosts at every large scale of
a set with long windows, with constants that depend only on $\varepsilon$. The definition allows an
extra requirement on each host, which for a primitive character $\psi$ is the host conditions of its
cyclic level and a bound on Tate-Shafarevich groups.

% before: hf-hosts-on-windows
For a primitive character $\psi$ that satisfies the Tate-Shafarevich bound, the host fields built
from $\psi$, with good pairs chosen on the windows where that bound holds, extend to good hosts.
These are the hosts that Chapter~\ref{chapter-assembly} requires for the cyclic level of $\psi$.
